\section{The Shifted Multi-Preimage Trapdoor Sampler Construction.}
We recall the modified shifted multi-preimage trapdoor sampler described in ~\cite[Construction~B.3]{EPRINT:BonLauTas25}.
\label{ap:1}
Let
$
    n=n(\lambda,\ell,d),
    q=q(\lambda,\ell,d),
    \sigma=\sigma(\lambda,\ell,d),
$
and set
$
    m = 3n\lceil \log q\rceil,
    t = \tilde{m}(d+1).
$
The vectors $u_i$ are drawn from
$\
    \{0,1\}^{d}.
$

\begin{itemize}
    \item 
    $\mathsf{Gen}(1^\lambda,1^\ell)$ is a randomized algorithm that samples
    $
        [\mathbf A |\mathbf B] \leftarrow \mathbb Z_q^{n\times t},
    $
    where
    $
        \mathbf A\in \mathbb Z_q^{n\times m},
        \mathbf B\in \mathbb Z_q^{n\times \tilde{m}d},
    $
    and outputs
    $
        \mathsf{crs}:=[\mathbf A| \mathbf B].
    $

    \item
    $\mathsf{Expand}(1^\lambda,1^\ell,1^d,\mathsf{crs},(u_i)_{i\in[\ell]})$:
    Given
    $
        \mathsf{crs}:=[\mathbf A\mid \mathbf B],
    $
    the deterministic expansion algorithm proceeds as follows.
    \begin{itemize}
        \item For each $i\in[\ell]$, define
    $
        \mathbf B_i := \mathbf B - u_i^\top \otimes \mathbf G,
    $
    and
    \[
        \mathbf A_i
        :=
        [\mathbf A \mid \mathbf B-u_i^\top\otimes \mathbf G]
        =
        [\mathbf A\mid \mathbf B_i]
        \in \mathbb Z_q^{n\times t}.
    \]

    \item Let
    $
        \Pi\in \{0,1\}^{(\ell t+\tilde{m})\times(\ell t+\tilde{m})}
    $
    be the permutation matrix such that
    \[
\mathbf{D}_\ell
:=
\left[
\begin{array}{ccccc|c}
    \mathbf A & \mathbf B_1 &             &             &              & \mathbf G \\
              & \mathbf A   & \mathbf B_2 &             &              & \mathbf G \\
              &             & \ddots      & \ddots      &              & \vdots    \\
              &             &             & \mathbf A   & \mathbf B_\ell & \mathbf G
\end{array}
\right]
=
\left[
\begin{array}{ccc|ccc|c}
    \mathbf A &        &           & \mathbf B_1 &        &              & \mathbf G \\
              & \ddots &           &              & \ddots &              & \vdots    \\
              &        & \mathbf A &              &        & \mathbf B_\ell & \mathbf G
\end{array}
\right]\Pi
\in \mathbb Z_q^{\ell n \times (\ell t+\tilde{m})} .
\]

    \item Compute
    $
        \mathbf T'
        \leftarrow
        \mathsf{StructTrapGen}(\mathbf B,(u_i)_{i\in[\ell]})
        \in
        \mathbb Z_q^{(\ell \tilde{m}d+\tilde{m})\times \ell \tilde{m}},
    $
    where $\mathsf{StructTrapGen}$ is the algorithm from
    \cite[Lemma~4.1]{EC:WatWeeWu25}. It outputs a trapdoor $\mathbf T'$
    for the matrix
    \[
\left[
\begin{array}{cccc|c}
    \mathbf B - u_1^\top \otimes \mathbf G
        &        &        &        & \mathbf G \\
    & \mathbf B - u_2^\top \otimes \mathbf G
        &        &        & \mathbf G \\
    &        & \ddots &        & \vdots \\
    &        &        &
    \mathbf B - u_\ell^\top \otimes \mathbf G
        & \mathbf G
\end{array}
\right].
\]
    such that
    $
        \|\mathbf T'\|_\infty = 1.
    $

    \item Finally, the procedure defines
    $
        \mathbf T\in \mathbb Z_q^{(\ell t+\tilde{m})\times \ell \tilde{m}}
    $
    as
    \[
        \mathbf T
        :=
        \Pi^{-1}
        \begin{bmatrix}
            \mathbf 0_{\ell \tilde{m}\times \ell \tilde{m}}\\
            \mathbf T'
        \end{bmatrix}.
    \]
    The algorithm sets
    $
        \mathsf{td}:=(\mathbf D_\ell,\mathbf T)
    $
    and outputs
    $
        (\mathbf A_1,\ldots,\mathbf A_\ell,\mathsf{td}).
    $
\end{itemize}
    \item
    $\mathsf{SampleMultPre}(\mathsf{td},\mathbf t_1,\ldots,\mathbf t_\ell)$:
    Given the trapdoor
    $
        \mathsf{td}:=(\mathbf D_\ell,\mathbf T)
    $
    and target vectors
    $
        \mathbf t_1,\ldots,\mathbf t_\ell\in\mathbb Z_q^n,
    $
    define
    $
        \widehat{\mathbf t}
        :=
        \mathsf{col}(\mathbf t_1,\ldots,\mathbf t_\ell).
    $
    Then sample
    \[
        \mathsf{col}(\boldsymbol\pi_1,\ldots,\boldsymbol\pi_\ell,\widehat{\mathbf c})
        \leftarrow
        \mathsf{SamplePre}
        (\mathbf D_\ell,\mathbf T,\widehat{\mathbf t},\sigma),
    \]
    and output
    $
        (\boldsymbol\pi_1,\ldots,\boldsymbol\pi_\ell,-\mathbf G\widehat{\mathbf c}).
    $
    
    \item
    $\mathsf{GenProg}(1^\lambda,1^\ell,1^d,u_i,\mathbf A_i)$:
    This deterministic algorithm takes a vector
    $
        u_i\in\{0,1\}^{\lceil\log\ell\rceil}
    $
    and a matrix
    $
        \mathbf A_i:=[\mathbf A\mid \mathbf B],
    $
    and outputs
    $
        \mathsf{crs}:=
        [\mathbf A\mid \mathbf B+u_i^\top\otimes \mathbf G].
    $

    \item
    $\mathsf{ExpandLocal}(1^\lambda,\mathsf{crs},u_i)$:
    This deterministic algorithm takes
    $
        \mathsf{crs}:=[\mathbf A\mid \mathbf B]
    $
    and
    $
        u_i\in\{0,1\}^{d},
    $
    and outputs
    $
        \mathbf A_i
        :=
        [\mathbf A\mid \mathbf B-u_i^\top\otimes \mathbf G]
        \in \mathbb Z_q^{n\times t}.
    $
\end{itemize}
\section{The construction of SampleLeft}
Let $\mathbf A \in \mathbb Z_q^{n \times m}$, 
$\mathbf M_1 \in \mathbb Z_q^{n \times m_1}$, and let $\sigma$ be a Gaussian parameter. In essence, to find the short preimage of $\mathbf{F}=[\mathbf{A}|\mathbf{M}_1]$, the $\mathsf{SampleLeft}$ employs the basis extension method discovered in~\cite{EC:CasHofKilPei10},which constructs a basis $\mathbf{T_F}$ of the orthogonal lattice of basis $\mathbf{F}$, such that $\|\mathbf{T_F}\|_{\mathsf{GS}}=\|\mathbf{T_A}\|_{\mathsf{GS}}$. After that, it simply invokes the algorithm $\mathsf{SamplePre}(\mathbf{F},\mathbf{T_F},\mathbf{u},\sigma)$ to obtain the desirable vector $\mathbf{v}$. 

Suppose $\mathbf{T_A} \in \mathbb{Z}_q^{m \times m}$ is  a corresponding Aitaj trapdoor of $\mathbf{A}$, there is a deterministic algorithm $\mathsf{DetExtBasis}(\mathbf{T_A},[\mathbf{A}|\mathbf{M}_1])$ described below:

Let
$
    \mathbf{T}
    :=
    \begin{bmatrix}
        \mathbf{t}_1 & \cdots & \mathbf{t}_{m_1}
    \end{bmatrix}
    \in \mathbb{Z}^{m \times m_1}
$
be a matrix whose columns satisfy
$
    \mathbf{A}\mathbf{t}_i = -\mathbf{M}_1[i] \pmod q,
    i \in [m_1].
$, where $t_i$ is found by using standard Gaussian elimination.
Equivalently,
$
    \mathbf{A}\mathbf{T} = -\mathbf{M}_1 \pmod q .
$. We then define
\[
    \mathbf{T_F}
    :=
    \begin{bmatrix}
        \mathbf{T_A} & \mathbf{T} \\
        \mathbf{0}_{m_1 \times m} & \mathbf{I}_{m_1}
    \end{bmatrix}
    \in \mathbb{Z}^{(m+m_1)\times(m+m_1)} .
\]
, which is a valid trapdoor of $\Lambda^{\perp}(\mathbf{F})$



By calling $\mathsf{SamplePre}(\mathbf{F},\mathbf{T_F},\mathbf{u},\sigma)$ after computing $\mathbf{T_F}$ as described above, we can derive a preimage $\mathbf{u} \in \mathbb{Z}_q^{m + m_1}$ of syndrome vector $\mathbf{u} \in \mathbb{Z}_q^n$, which has distribution statistically close to $\mathbf{F}^{-1}_\sigma(\mathbf{v})$

Interestingly, this algorithm is not restricted to Ajtai trapdoors. With a few additional steps, it can also be adapted to take a gadget trapdoor as input. To justify this observation, we first recall several useful lemmas that allow gadget trapdoors to be used within $\mathsf{SampleLeft}$.
\begin{lemma}[Gadget trapdoor yields an Ajtai trapdoor~\cite{EPRINT:ChaHsiWu25}]
\label{lemA:1}
Let $n,q$ be lattice parameters and let $m' = n\lceil \log q\rceil$.
For any matrix $\mathbf A \in \mathbb Z_q^{n \times m}$ and any
$\mathbf T \in \mathbb Z_q^{m \times m'}$ satisfying
$
    \mathbf A \mathbf T = \mathbf G_n \pmod q,
$
let
$
    \mathbf T'
    =
    \left[
        \mathbf I_m - \mathbf T \mathbf G_n^{-1}(\mathbf A)
        \ \middle|\ 
        \mathbf T \mathbf S_n
    \right]
    \in \mathbb Z_q^{m \times (m+m')},
$
where $\mathbf S_n \in \mathbb Z_q^{m' \times m'}$ is an Ajtai trapdoor
for the gadget matrix $\mathbf G_n$.
Then $\mathbf T'$ is an Ajtai trapdoor for $\mathbf A$; in particular,
$
    \mathbf A \mathbf T' = \mathbf 0 \pmod q,
$
and $\mathbf T'$ has full rank over $\mathbb R$.
\end{lemma}

\begin{lemma}[SampleLeft with a gadget trapdoor]
\label{ap:2}
Let $q \geq 2$, $m > n$, and let $m' = n\lceil \log q\rceil$.
Let
$
    \mathbf A \in \mathbb Z_q^{n \times m},
    \mathbf M_1 \in \mathbb Z_q^{n \times m_1},
$
and define
$
    \mathbf F = [\mathbf A | \mathbf M_1]
    \in \mathbb Z_q^{n \times (m+m_1)}.
$
Assume that $\widehat{\mathbf T}_{\mathbf A} \in \mathbb Z_q^{m \times m'}$
is a gadget trapdoor for $\mathbf A$, namely
$
    \mathbf A\widehat{\mathbf T}_{\mathbf A}
    =
    \mathbf G_n
    \pmod q .
$
Then, for any $\mathbf u \in \mathbb Z_q^n$ and any Gaussian parameter
$
    \sigma
    >
    2m^{3/2}\|\widehat{\mathbf T}_{\mathbf A}\|_{\infty}
    \cdot \omega(\log(m+m_1)),
$
there exists an efficient probabilistic algorithm
$
    \mathsf{SampleLeft}_{\mathsf{G}}
    (\mathbf A,\mathbf M_1,\widehat{\mathbf T}_{\mathbf A},\mathbf u,\sigma)
$
that outputs a vector
$
    \mathbf v \in \Lambda_q^{\mathbf u}(\mathbf F)
$
distributed statistically close to
$
    \mathbf{F}_\sigma^{-1}(\mathbf{u)}
$
\end{lemma}
\begin{proof}
We first invoke Lemma~\ref{lemA:1} to convert the gadget trapdoor into an
Ajtai-style trapdoor and then extract $m$ first column from its. This yields
$
    \mathbf{T}_{\mathbf A} \in \mathbb{Z}_q^{m \times m}.
$
We then sample a preimage of $\mathbf{u}$ using the Ajtai-trapdoor version of
$\mathsf{SampleLeft}$. By Eq.~(7.1) in Lemma~7.8 of~\cite{EPRINT:ChaHsiWu25}, the Gram--Schmidt norm
of $\mathbf{T}_{\mathbf A}$ satisfies
$
    \left\lVert \mathbf{T}_{\mathbf A} \right\rVert_{\mathsf{GS}}
    \leq
    \sqrt{m}
    \left\lVert \mathbf{T}_{\mathbf A} \right\rVert_{\infty}
    \leq
    2m^{3/2}
    \| \widehat{\mathbf{T}}_{\mathbf A} \|_{\infty}.
$
Therefore, by Lemma~\ref{lem:gadget_trapdoor}, it suffices to choose the Gaussian parameter
$ \sigma
    >
    2m^{3/2}
    \| \widehat{\mathbf{T}}_{\mathbf A} \|_{\infty}
    \cdot
    \omega\!\left(\log(m+m_1)\right).
$
This completes the proof.
\end{proof}

\section{The construction of Explainable SampleLeft}
\label{ap:3}

Let $\mathbf{A}\in\mathbb{Z}_q^{n\times m}$,
$\mathbf{M}_1\in\mathbb{Z}_q^{n\times m_1}$, and define
$
    \mathbf{F}:=[\mathbf{A}\mid \mathbf{M}_1]
    \in \mathbb{Z}_q^{n\times(m+m_1)} .
$

\begin{itemize}[leftmargin=*]

    \item 
    $\mathsf{SampleLeft}(\mathbf{A},\mathbf{M}_1,\mathbf{T}_{\mathbf A},
    \mathbf{u},\sigma;\mathsf{rb})$ is probabilistic algorithm that  as follows

    \begin{itemize}
        \item First, convert the gadget trapdoor $\mathbf{T}_{\mathbf A}$ into an
        Ajtai-style trapdoor. Let $\mathbf{T}'_{\mathbf A}$ be the first $m$
        linearly independent columns over $\mathbb{R}$ of
        \[
            \left[
                \mathbf{I}_m
                -
                \mathbf{T}_{\mathbf A}\mathbf{G}^{-1}\mathbf{A}
                \;\middle|\;
                \mathbf{T}_{\mathbf A}\mathbf{S}
            \right],
        \]
        where $\mathbf{S}\in\mathbb{Z}^{m\times m}$ is an Ajtai trapdoor for
        $\mathbf{G}$. Then $\mathbf{T}'_{\mathbf A}$ is an Ajtai trapdoor for
        $\mathbf{A}$.

        \item Next, deterministically extend this trapdoor to the matrix
        $\mathbf{F}=[\mathbf{A}\mid\mathbf{M}_1]$ by computing
        $
            \mathbf{T}'_{\mathbf F}
            :=
            \mathsf{DetExtBasis}
            \bigl(
                \mathbf{T}'_{\mathbf A},
                [\mathbf{A}\mid\mathbf{M}_1]
            \bigr).
        $

        \item Finally, output
        $
            \mathbf{x}
            \leftarrow
            \mathsf{SamplePre}
            \bigl(
                \mathbf{F},
                \mathbf{T}'_{\mathbf F},
                \mathbf{u},
                \sigma;
                \mathbf{r_b}
            \bigr).
        $
    \end{itemize}

    \item
    $\mathsf{ExplainSL}(1^\lambda,\mathbf{A},\mathbf{M}_1,
    \mathbf{T}_{\mathbf A},\mathbf{u},\mathbf{y},\sigma)$ is a deterministic algorithm, which proceeds as follows.

    \begin{itemize}
        \item Compute the extended trapdoor
        $
            \mathbf{T}_{\mathbf F}
            :=
            \mathsf{DetExtBasis}
            \bigl(
                \mathbf{T}_{\mathbf A},
                [\mathbf{A}\mid\mathbf{M}_1]
            \bigr).
        $

        \item Then output
        $\mathsf{rb} :=
            \mathsf{ExplainSP}
            \bigl(
                1^\lambda,
                \mathbf{F},
                \mathbf{T}_{\mathbf F},
                \mathbf{u},
                \mathbf{y},
                \sigma
            \bigr),
        $
        where the Ajtai trapdoor used in $\mathsf{ExplainSP}$ is replaced by
        $\mathbf{T}_{\mathbf F}$.
    \end{itemize}

\end{itemize}

\section{Partial gadget trapdoor}
\label{ap:4}
Let $k,N,n,\log q \in poly(\lambda)$ and
$m = n \cdot (\lceil \log q \rceil + 1)$.
Let $\{v_0,v_1,\ldots,v_{N-1}\}\subseteq \mathbb{Z}_q^{\times}$
be an arbitrary set such that
$v_i - v_j \in \mathbb{Z}_q^{\times}$ for all $i\neq j$.
For $j\in[N]$, let
$
    \mathbf{v}_j := (1,v_j,\ldots,v_j^{k-1})^{\top},
$
$\widehat{\mathbf{e}}_{T,j}\in\{0,1\}^{T}$ be the unit vector indexed by $j$,
and $(\mathsf{TrapGen},\mathsf{SampPre})$ described in Lemma~\ref{lem:gadget_trapdoor} be a lattice trapdoor scheme,
which is $(nk,2mk,\widetilde{s})$-correct with some $\widetilde{s}>0$.
Finally, $\mathbf{S}\in R^{2mk\times 2mk}$ is the Gaussian parameter
hardwired in the $\mathsf{PTrapGen}$ algorithm.

\begin{itemize}
    \item $\mathsf{PTrapGen}(1^\lambda,\mathbf{A},\mathsf{td},j)$ is a randomized
    algorithm that generates the partial trapdoor for party $j$. It proceeds as
    follows.

    \begin{itemize}
        \item First, sample an auxiliary trapdoored matrix
        $
            (\mathbf{C}_j,\mathsf{td}_{\mathbf{C}_j})
            \gets
            \mathsf{TrapGen}(1^\lambda,1^n,1^m),
        $
        where $\mathbf{C}_j\in R_q^{n\times m}$.

        \item Then, using the master trapdoor $\mathsf{td}$ of $\mathbf{A}$,
        sample
        $
            \mathbf{U}_j
            \gets
            \mathsf{SamplePre}
            (\mathbf{A},\mathsf{td},
            \mathbf{v}_j\otimes \mathbf{C}_j,\mathbf{S}).
        $
        Thus, $\mathbf{U}_j$ satisfies
        $
            \mathbf{A}\mathbf{U}_j
            \equiv
            \mathbf{v}_j\otimes \mathbf{C}_j
            \pmod q.
        $

        \item If $\operatorname{Rank}_{K}(\mathbf{U}_j)\neq m$, the algorithm
        outputs $\bot$.

        \item Otherwise, it outputs the partial trapdoor
        $
            \mathsf{ptd}_j
            :=
            \left(
                \mathbf{U}_j,
                \mathbf{C}_j,
                \mathbf{T}_{\mathbf{C}_j}
            \right).
        $
    \end{itemize}

    \item $\mathsf{PSampPre}(\mathbf{A},\mathsf{ptd}_j,T,\mathbf{y},\sigma)$
    is a randomized algorithm that samples a partial preimage for party $j$
    with respect to the reconstruction set $T$. It proceeds as follows.

    \begin{itemize}
        \item Let
        $
            \mathbf{V}_T := (\mathbf{v}_j)_{j\in T}
            \in \mathbb{Z}_q^{k\times k}.
        $

        \item Decompose the target image $\mathbf{y}$ according to the basis
        indexed by $T$. Namely, compute
        $
            \mathbf{z}_j
            :=
            \left(
                \widehat{\mathbf{e}}_{T,j}^{\top}
                \mathbf{V}_T^{-1}
                \otimes \mathbf{I}_n
            \right)\mathbf{y}
            \pmod q.
        $
        This gives the unique decomposition
        \[
            \mathbf{y}
            \equiv
            \sum_{j\in T}\mathbf{v}_j\otimes \mathbf{z}_j
            \pmod q.
        \]

        \item Define the Gaussian covariance parameter
        $
            \boldsymbol{\Sigma}_j
            :=
            \sigma^2
            (\mathbf{U}_j^{\top}\mathbf{U}_j)^{-1}.
        $

        \item Using the trapdoor $\mathsf{td}_{\mathbf{C}_j}$ of
        $\mathbf{C}_j$, sample
        $
            \mathbf{d}_j
            \gets
            \mathsf{SamplePre}
            (\mathbf{C}_j,\mathbf{T}_{\mathbf{C}_j},
            \mathbf{z}_j,\sqrt{\boldsymbol{\Sigma}_j}),
        $
        so that
        $
            \mathbf{C}_j\mathbf{d}_j
            \equiv
            \mathbf{z}_j
            \pmod q.
        $
        \item Finally, output
        $
            \mathbf{x}_j
            :=
            \mathbf{U}_j\mathbf{d}_j
            \in \mathbb{Z}^{2mk}.
        $
        By construction,
        \[
            \mathbf{A}\mathbf{x}_j
            =
            \mathbf{A}\mathbf{U}_j\mathbf{d}_j
            \equiv
            (\mathbf{v}_j\otimes \mathbf{C}_j)\mathbf{d}_j
            \equiv
            \mathbf{v}_j\otimes \mathbf{z}_j
            \pmod q.
        \]
    \end{itemize}

    \item $\mathsf{Rec}((\mathbf{x}_j)_{j\in T})$ is a deterministic algorithm
    that reconstructs a full preimage from the partial preimages. It simply
    computes
    $
        \mathbf{x}
        :=
        \sum_{j\in T}\mathbf{x}_j.
    $
    Then,
    \[
        \mathbf{A}\mathbf{x}
        =
        \sum_{j\in T}\mathbf{A}\mathbf{x}_j
        \equiv
        \sum_{j\in T}\mathbf{v}_j\otimes \mathbf{z}_j
        \equiv
        \mathbf{y}
        \pmod q.
    \]
    Hence, $\mathbf{x}$ is a valid full preimage of $\mathbf{y}$ under
    $\mathbf{A}$.
\end{itemize}